% Overleaf: select pdfLaTeX, then Recompile.
% For Chinese text: select XeLaTeX and uncomment the ctex line below.
\documentclass[a4paper,11pt]{article}
% \usepackage[UTF8,scheme=plain]{ctex}
\usepackage{hku-cds-assignment}

% ---------- Edit your assignment details here ----------
\newcommand{\CourseCode}{COMPXXXX}
\newcommand{\CourseName}{Course Name}
\newcommand{\AssignmentName}{Assignment 1}
\newcommand{\AssignmentTitle}{Assignment Title}
\newcommand{\StudentName}{Your Name}
\newcommand{\StudentID}{303XXXXXXXX}
\newcommand{\SubmissionDate}{6 November 2026}

\coverpagetrue % Change to \coverpagefalse for a compact first-page heading.

\begin{document}
\makeassignmenttitle

% Replace all example answers below with your own work.
\question{Mathematical analysis}
\subsection*{(a) Problem formulation}
Define the task, assumptions and notation before presenting the solution.
For example, let $X\in\mathbb{R}^{n\times d}$ be a design matrix and
$\mathbf{y}\in\mathbb{R}^{n}$ the response vector. Consider the ridge
regression objective
\begin{equation}
  \mathcal{L}(\mathbf{w}) = \frac{1}{n}\norm{X\mathbf{w}-\mathbf{y}}_2^2
  + \lambda\norm{\mathbf{w}}_2^2, \qquad \lambda>0.
  \label{eq:objective}
\end{equation}

\subsection*{(b) Solution}
Differentiating Equation~\eqref{eq:objective} and setting the gradient to
zero gives
\begin{align}
  \nabla_{\mathbf{w}}\mathcal{L}
    &= \frac{2}{n}X^{\mathsf T}(X\mathbf{w}-\mathbf{y})+2\lambda\mathbf{w},\\
  \widehat{\mathbf{w}}
    &= (X^{\mathsf T}X+n\lambda I_d)^{-1}X^{\mathsf T}\mathbf{y}.
\end{align}
The positive regularization parameter makes the system matrix positive
definite. In code, solve the linear system directly rather than explicitly
computing its inverse.

\question{Implementation and results}
\subsection*{(a) Method}
Describe the data, preprocessing, evaluation metric and experimental setup.
Use a fixed seed where appropriate, and fit preprocessing steps on training
data only. Keep enough detail for the experiment to be reproduced.

\begin{lstlisting}[language=Python,float=htbp,caption={A minimal ridge regression implementation.},label={lst:ridge}]
import numpy as np

def fit_ridge(X, y, lam=0.1):
    n, d = X.shape
    A = X.T @ X + n * lam * np.eye(d)
    return np.linalg.solve(A, X.T @ y)
\end{lstlisting}

\ifcoverpage\clearpage\fi % Optional page break in the cover-page version.
\subsection*{(b) Results and discussion}
Table~\ref{tab:results} illustrates a simple results table. The values are
fictional examples; replace them with measured results and explain the
observed differences.

\begin{table}[htbp]
  \centering
  \caption{Illustrative results only (lower RMSE is better).}
  \label{tab:results}
  \begin{tabular}{lrr}
    \toprule
    Method & Validation RMSE & Test RMSE \\
    \midrule
    Mean baseline & 1.24 & 1.28 \\
    Ridge regression & 0.87 & 0.91 \\
    \bottomrule
  \end{tabular}
\end{table}

\subsection*{(c) Figures}
Replace the box below with your own figure. For example:
\begin{quote}
\small\verb|\includegraphics[width=.75\linewidth]{your-figure.pdf}|
\end{quote}

\begin{figure}[htbp]
  \centering
  \fbox{\parbox[c][30mm][c]{0.72\linewidth}{\centering
    Your figure here\\[4pt]
    \small Plot, diagram, or experimental output}}
  \caption{Replace this placeholder and caption with your own figure.}
  \label{fig:example}
\end{figure}

\question{Conclusion}
Summarize the main finding, explain one limitation, and relate the result to
the question. Cite external sources when appropriate, for example the
reference repository used for the visual inspiration of this template~\cite{template}.

% For a short assignment, thebibliography avoids a separate BibTeX step.
\begin{thebibliography}{9}
\raggedright
\bibitem{template}
Guanying Chen.
\emph{HKU-PhD-Thesis-LaTex}.
\url{https://github.com/guanyingc/HKU-PhD-Thesis-LaTex}.
\end{thebibliography}

% Optional appendix:
% \appendix
% \section{Additional details}

\end{document}
